Theorem~\ref{thm:limite} already implies that $\lim_nf_h(n)$ is positive for each fixed $h$. Indeed, \eqref{eq:F} and the bounds in
Theorem~\ref{thm:perfil}(a) and~(c) give $F(0)\ge\kappa^2>a_0/2$. The direct proof of Theorem~\ref{thm:positiva} is in
Section~\ref{sec:positiva}. The constant is conservative, and the proof gives only the existence of the threshold $n_0$.

To understand this proof, return to the case of the last queen. If the crossing $(r,c)$ is under attack, direct insertion there fails.
After removing a queen from $(i,j)$, we can occupy its \emph{destinations} $(r,j)$ and $(i,c)$ if the two new queens are compatible both
with the rest of the board and with each other. This fills row $r$ and column $c$ even though the crossing $(r,c)$ remains unoccupied.

For any $h\ge1$, we call this replacement of one queen by two a \emph{one-move repair} at $(r,c)$. A legal repair takes the placement
$P\in\Om_h(n)$ to layer $h-1$. Figure~\ref{fig:reparacion-tablero} shows an example on the $4\times4$ board. The queen at $(3,2)$ attacks the
empty crossing $(2,1)$. We remove the queen at $(0,0)$ and place two queens at $(2,0)$ and $(0,1)$. The original crossing remains under
attack and unoccupied. The resulting placement has four queens and is the one in Figure~\ref{fig:ejemplos}(a).

\begin{figure}[htbp]
\centering
\begin{tikzpicture}[tablero A]
 \begin{scope}
  \cuadriculaA{4}\trazarRejillaA{4}
  \ataqueA{3}{2}{2}{1}\aspaA{2}{1}
  \aroA{0}{0}\reinaA[colorAro]{0}{0}\reinaA{1}{3}\reinaA{3}{2}
  \panelA{2,4.38}{(a) Remove one queen}
 \end{scope}
 \draw[-{Stealth[length=2mm]},tintaReina,line width=.9pt] (4.9,2)--(6.65,2);
 \begin{scope}[shift={(8,0)}]
  \cuadriculaA{4}\trazarRejillaA{4}
  \ataqueA{3}{2}{2}{1}\aspaA{2}{1}
  \reinaA{1}{3}\reinaA{3}{2}
  \aroA{0}{1}\reinaA[colorAro]{0}{1}
  \aroA{2}{0}\reinaA[colorAro]{2}{0}
  \panelA{2,4.38}{(b) Place two queens}
 \end{scope}
\end{tikzpicture}
\caption{A legal repair takes a placement from $\Omega_1(4)$ to $\Omega_0(4)$.
The rings mark the pieces involved. The cross marks the same square $(2,1)$ in both panels.
Reading from right to left recovers the reverse operation with the ordered pair $(2,0),(0,1)$.}
\label{fig:reparacion-tablero}
\end{figure}


We call a queen that can legally be replaced in this way an \emph{absorber} for the crossing. In
Figure~\ref{fig:reparacion-tablero}(a), only the queen at $(0,0)$ qualifies. If we remove the queen at $(1,3)$, one of its destinations is
$(1,1)$. If we remove the queen at $(3,2)$, one of its destinations is $(2,2)$. In both cases, the queen at $(0,0)$ remains in place and
attacks that destination. Thus this crossing admits no direct insertion and exactly one legal repair.

Let $A^{rc}(P)$ be the number of legal repairs at $(r,c)$ and let $d_{h,1}(P)=\sum A^{rc}(P)$, where the sum runs over the $h^2$ empty
crossings. The subscript $1$ refers to one-move repairs, in which a single queen is removed. We write $\bar r=(r+\frac12)/n$ and $\bar
c=(c+\frac12)/n$ for the coordinates of the centers of the squares. To express how the number of repairs varies with the position of the
crossing, we use the full-board profile. We denote its density by $\nu_1$, whose subscript is the mass of queens, equal to one,
corresponding to $x=0$. We drop the subscript $x$ from its factors (the constant retains its subscript, $b_0$) and set \begin{gather*}
 \nu_1=b_0\,r_R\,r_C\,u_D\,u_A,\qquad \psi(r,c)=u_D(r-c)\,u_A(r+c),\\
 a(r,c)=\int_{[0,1]^2}\nu_1(i,j)\,\psi(r,j)\,\psi(i,c)\,di\,dj .
\end{gather*} Then $\kappa^4\le a\le1$.

\begin{teorema}[Typical structure]\label{thm:estructura} Fix $h\ge1$ and let $P$ have law $\mu_h$. \begin{enumerate}[(a)] \item
\emph{One-move repairs.} In probability, \[
 \max_{(r,c)\ \text{empty crossing of }P}\ \Bigl|\frac{A^{rc}(P)}{n}-a(\bar r,\bar c)\Bigr|\longrightarrow0 .
\] \item \emph{Mean spatial distribution of the holes.} For every continuous function $t$ on $[0,1]^2$, \[
 \E_{\mu_h}\Bigl[\frac1{h^2}\sum_{(r,c)\ \text{empty crossing}}t(\bar r,\bar c)\Bigr]
 \longrightarrow\int_{[0,1]^2}t(r,c)\,r_R(r)\,r_C(c)\,dr\,dc .
\] For $h=1$ this gives weak convergence of the law of the position of the unique empty crossing. \item \emph{Unweighted counter.}
$\E_{\mu_h}[d_{h,1}]/(h^2n)\to B_1/b_0$, where $B_1=\int_{[0,1]^2}\psi$. \end{enumerate} \end{teorema}

\begin{observacion}[Scope]\label{obs:alcance} \begin{enumerate}[(i)] \item All statements hold for fixed $h$ and $K$ as $n\to\infty$.
We obtain neither convergence rates nor uniformity for $K=K(n)\to\infty$, and we do not determine the full asymptotics of $Q(n)$. \item
For $h>1$, part (b) identifies only the mean measure of the empty crossings. It makes no claim about asymptotic independence of the
empty rows and columns or concentration of the empirical measure of the empty crossings. Part (c) concerns convergence of means. By
(a), $d_{h,1}/n-\sum a(\bar r,\bar c)\to0$, and the second term may remain random, so we make no claim about concentration of
$d_{h,1}/n$. \item The value $F(0)\approx0.0809$ comes from a deterministic numerical solution of the variational
problem~\cite[Section~9]{Perfil}, with no certified error bound. At finite board sizes, the exact values from exhaustive enumeration
are far from the limit (for example $f_1(15)\approx0.114$ and $f_2(15)\approx0.153$, see Table~\ref{tab:ilustracion}). The theorem
gives no information on the rate of convergence. \end{enumerate} \end{observacion}

\subsection{Outline of the proof}\label{subsec:idea}

Figure~\ref{fig:esquema} summarizes the structure of the proof.

A placement may admit several repairs at a single crossing. We count each legal operation once from its initial placement and once from
its final placement. In the second description we start from $P'\in\Om_{h-1}$, choose the ordered pair of queens $(r,j),(i,c)$, and
replace them by $(i,j)$ to recover $P\in\Om_h$. We assign the crossing $(r,c)$ the same weight $w(r,c)$ in both descriptions. The
double count in Section~\ref{sec:doblecuenta} gives \[
 N_h\,\E_{\mu_h}D_w=N_{h-1}\,\E_{\mu_{h-1}}I_w,\qquad\text{hence}\qquad
 f_h=\frac{n-h+1}{n}\,\frac{\E_{\mu_h}[D_w/(h^2n)]}{\E_{\mu_{h-1}}[I_w/n^2]},
\] where $D_w$ and $I_w$ are the weighted forward and reverse counters. To use the ratio, we will choose a weight for which the
denominator has a positive limit.

The two means are taken under different laws, $\mu_h$ for the forward count and $\mu_{h-1}$ for the reverse count. To handle both,
Section~\ref{sec:referencia} constructs a random process for each layer $h$. It places real pieces (queens) and virtual pieces
successively until no active line remains. A virtual piece places no queen and deactivates the lines it touches. The weights come from
the variational profile. Let $\nu'_h$ denote the law of the final placement in layer $h$, conditional on the success of the process. An
entropy identity, estimates for the normalizers (the sums of weights that normalize each choice), and Simkin's first-order lower bound
give $\KL{\mu_h}{\nu'_h}=o(n)$. Consequently, an event with probability at most $e^{-cn}$ under $\nu'_h$, for fixed $c>0$, has
probability $o(1)$ under $\mu_h$.
